\documentclass[10pt]{article}
\usepackage[T1]{fontenc}
\usepackage[utf8]{inputenc}
\usepackage[french]{babel}
\usepackage{amsmath,amssymb,mathrsfs,mathpazo}
\usepackage{geometry}
\geometry{a4paper,margin=2cm}
\usepackage{graphicx,array,tabularx,tabularray,colortbl,multirow,multicol,booktabs,enumitem,pifont,ulem,fancybox,diagbox,cancel}
\usepackage{pstricks,pstricks-add,pst-plot,pst-tree,pst-node,pst-text,pst-func,pst-math,pst-bspline,pst-3dplot,pst-eucl,pst-solides3d}
\let\pscalculate\undefined
\usepackage{pst-bezier}
\usepackage{tikz,pgfplots,tkz-tab}
\usetikzlibrary{arrows,arrows.meta,patterns,calc,positioning,trees,decorations.markings,shapes.geometric,shapes.misc,shapes.arrows,intersections,decorations.pathmorphing,decorations.pathreplacing,backgrounds,shadows,babel}
\pgfplotsset{compat=1.18}
\usepgfplotslibrary{fillbetween}
\usepackage[np]{numprint}
\usepackage{listings,makecell,pas-tableur,twoopt}
\let\barre\undefined
\usepackage[pstricks]{tablvar}
\definecolor{darkgreen}{RGB}{0,100,0}
\providecommand{\red}{\color{red}}
\providecommand{\blue}{\color{blue}}
\providecommand{\green}{\color{green}}
\providecommand{\black}{\color{black}}
\providecommand{\white}{\color{white}}
\providecommand{\gray}{\color{gray}}
\providecommand{\cyan}{\color{cyan}}
\providecommand{\magenta}{\color{magenta}}
\providecommand{\yellow}{\color{yellow}}
\providecommand{\violet}{\color{violet}}
\providecommand{\orange}{\color{orange}}
\providecommand{\euro}{\texteuro}
\providecommand{\EUR}{\texteuro}
\providecommand{\ieme}{\textsuperscript{e}}
\providecommand{\degres}{\ensuremath{^\circ}}
\providecommand{\degree}{\ensuremath{^\circ}}
\providecommand{\vv}[1]{\overrightarrow{#1}}
\providecommand{\Coord}[1]{#1}
\providecommand{\decosix}{}
\providecommand{\starredbullet}{\ensuremath{\bullet}}
\pagestyle{empty}
\setlength{\parindent}{0pt}
\usepackage{adjustbox}\DeclareUnicodeCharacter{2236}{:}\DeclareUnicodeCharacter{2190}{\ensuremath{\leftarrow}}
\providecommand{\R}{\mathbb{R}}
\providecommand{\N}{\mathbb{N}}
\providecommand{\Z}{\mathbb{Z}}
\providecommand{\Q}{\mathbb{Q}}
\providecommand{\D}{\mathbb{D}}
\providecommand{\C}{\mathbb{C}}
\providecommand{\e}{{\mathrm{e}}}
\providecommand{\vect}[1]{\overrightarrow{#1}}
\providecommand{\vectt}[1]{\overrightarrow{\mathrm{#1}}}
\providecommand{\vvt}[1]{\overrightarrow{\mathrm{#1}}}
\providecommand{\np}[2][]{{#2~#1}}
\providecommand{\vv}[1]{{\overrightarrow{#1}}}
\renewcommand{\d}{{\mathrm{d}}}
\providecommand{\degres}{{^\circ}}
\providecommand{\degre}{{^\circ}}
\providecommand{\blue}{\color{blue}}
\providecommand{\red}{\color{red}}
\providecommand{\enskip}{\quad}
\providecommand{\llbracket}{\lbrack\!\lbrack}
\providecommand{\rrbracket}{\rbrack\!\rbrack}
\providecommand{\euro}{{\text{€}}}
\providecommand{\no}{{\text{n°}}}
\providecommand{\interval}[2]{{[#1;#2]}}
\providecommand{\intervFO}[2]{{[#1;#2[}}
\providecommand{\dotfill}{\ldots}

\providecommand{\ds}{\displaystyle}
\providecommand{\textding}[1]{#1}
\providecommand{\Oij}{\(\mathrm{O};\vec{\imath},\vec{\jmath}\)}
\providecommand{\Oijk}{\(\mathrm{O};\vec{\imath},\vec{\jmath},\vec{k}\)}
\providecommand{\Ouv}{\(\mathrm{O};\vec{u},\vec{v}\)}
\providecommand{\intervFF}[2]{\ensuremath{[#1;#2]}}
\renewcommand{\euro}{\ifmmode\text{\texteuro}\else\texteuro\fi}
\DeclareUnicodeCharacter{00B0}{\ifmmode^{\circ}\else\textdegree\fi}
\DeclareUnicodeCharacter{20AC}{\ifmmode\text{\texteuro}\else\texteuro\fi}
% Énoncés : Éducation nationale. Transcription et corrigés : APMEP, Denis Vergès et contributeurs. Découpage et classement : CoopMaths, Sébastien Lozano, Éric Elter et contributeurs. Adaptation de la présentation : Maths Explorer.
% CC BY-SA : https://coopmaths.fr/www/about/
% https://coopmaths.fr/alea/static/e3c/2020/tex/e3c_2020_00_sujet3_2_cor.tex
\begin{document}
\noindent\textit{APMEP / CoopMaths. CC BY-SA. Présentation adaptée par Maths Explorer.}\par\medskip
\section*{Exercice 2 (5 points)}
	

	
	\subsection*{1. Quelle est la valeur de $u_3$ ?}
	
	On a $u_3 = 18$.
	
	\subsection*{2. On admet que la suite $(u_n)$ est arithmétique de raison 6. Exprimer $u_n$ en fonction de $n$.}
	
	On a $u_n = 6n$.
	
	\subsection*{3. Combien l’artisan a-t-il ajouté de carreaux pour faire l’étape 5 ?}
	
	Avec $u_4 = 6 \times 4 = 24$ et $u_5 = 6 \times 5 = 30$, on a donc ajouté $30 - 24 = 6$ carreaux pour faire l’étape 5.
	
	\subsection*{4. Combien a-t-il alors posé de carreaux au total lorsqu’il termine l’étape 5 (en comptant le carreau central initial) ?}
	
	Il a posé en tout :
	\[
	1 + 6 + 12 + 18 + 24 + 30 = 91 \text{ carreaux}
	\]
	
	\subsection*{5. On pose $S_n = u_1 + u_2 + \dots + u_n$. Montrer que $S_n = 6(1 + 2 + 3 + \dots + n)$ puis que $S_n = 3n^2 + 3n$.}
	
	\[
	S_n = 6 \times (1 + 2 + 3 + 4 + 5 + \dots + n)
	\]
	
	Or
	\[
	T_n = 1 + 2 + 3 + \dots + n
	\]
	
	On peut écrire
	\[
	T_n = n + (n-1) + \dots + 3 + 2 + 1
	\]
	
	En sommant membre à membre :
	\[
	2T_n = (n+1) + (n+1) + \dots + (n+1) = n(n+1)
	\]
	
	Donc
	\[
	T_n = \dfrac{n(n+1)}{2}
	\]
	
	et par suite
	\[
	S_n = 6 \times \dfrac{n(n+1)}{2} = 3n(n+1) = 3n^2 + 3n
	\]
	
	\subsection*{6. Si on compte le premier carreau central, le nombre total de carreaux posés par l’artisan depuis le début, lorsqu’il termine la n-ième étape, est donc $3n^2 + 3n + 1$.}
	
	À la fin de sa semaine, l’artisan termine la pose du carrelage en collant son 2977ème carreau. Combien a-t-il fait d’étapes ?
	
	On a donc $3n^2 + 3n + 1 = 2977$. Il faut donc résoudre dans $\mathbb{N}$ l’équation :
	\[
	3n^2 + 3n - 2976 = 0
	\]
	
	ou en simplifiant par 3 :
	\[
	n^2 + n - 992 = 0
	\]
	
	On a
	\[
	\Delta = 1 + 4 \times 992 = 1 + 3968 = 3969 = 63^2
	\]
	
	Les solutions sont donc
	\[
	n_1 = \dfrac{-1 + 63}{2} = 31 \quad \text{et} \quad n_2 = \dfrac{-1 - 63}{2} = -32
	\]
	
	On ne retient que la solution 31. L’artisan a donc fait 31 étapes.

\end{document}
